\section{What the Comparison Yields}

The preceding sections were deliberately undramatic. This one states a judgement, gives the strongest case against it, and says what would overturn it. It is marked off from everything before it because the reader is entitled to know exactly where description stops.

\subsection{The two cases at full strength}

Before the judgement, both cases deserve their best form. Neither is given here in a polemicist's version; both are given in the words of authorities who legislated. That choice has a cost, stated here and not in the appendix alone: this study examined no work of the reform's non-official critical literature in an identified edition, so the case against appears below in official and papal formulations rather than in the words of the scholars who have pressed it at greatest length; a reader who thinks those scholars put it more strongly should treat what follows as a floor, not a ceiling.

\textbf{The case for the reform, at its strongest.} The reform's own argument is not that the older book was defective. It is that the norm invoked to justify the older book---restoration \emph{ad pristinam sanctorum Patrum normam}---is a norm the reform is entitled to invoke against it. The \emph{Institutio generalis} makes this argument in four articles and makes it carefully. Pius V's revisers, working in an age when the sacrificial character of the Mass, the ministerial priesthood and the real presence were under attack, rightly aimed to preserve ``the more recent tradition, unjustly assailed,'' introducing ``only the smallest changes'' to the sacred rite; and the 1570 Missal ``differs very little from the first printed edition of 1474, which in turn faithfully reproduces the Missal of the time of Innocent III.'' The manuscripts available to them did not permit an inquiry into ``the ancient and approved authors'' beyond the liturgical commentaries of the Middle Ages. Today that norm ``is enriched by innumerable writings of scholars''---the Gregorian Sacramentary printed in 1571, the old Roman and Ambrosian sacramentaries, the Spanish and Gallican books, the traditions of the centuries before East and West divided---and so ``the norm of the holy Fathers'' requires not only that what our nearer forebears handed down be kept, but that all past ages of the Church be comprehended and more deeply weighed.\endnote{\emph{Institutio Generalis Missalis Romani} (2002), \emph{Proœmium} nn.~6--9, artifact pp.~10--11. Number 7: ``Temporibus sane difficilibus, quibus catholica fides de indole sacrificali Missæ, de ministeriali sacerdotio, de reali et perpetua Christi sub eucharisticis speciebus præsentia in discrimen fuerat adducta, id S. Pii V imprimis intererat, ut recentiorem traditionem, immerito oppugnatam, servaret, minimis tantummodo ritus sacri mutationibus inductis.''} That is a serious historical argument, made from the sources, and it deserves to be met rather than dismissed: the older book is itself a moment in a tradition, shaped by a sixteenth-century emergency and by the limits of sixteenth-century scholarship, and fidelity to the Fathers may require going behind it.

\textbf{The case against the reform, at its strongest.} Its best form is not a claim of heresy; it is a claim about how a rite forms belief, and its most authoritative statement comes from a pope who legislated in the older book's favour. Benedict XVI, writing to the bishops of the world in 2007, granted that ``in many places celebrations were not faithful to the prescriptions of the new Missal,'' that ``the freedom to `creativity' \dots{} was often taken to extremes,'' and that this ``led to deformations of the liturgy which were hard to bear''; and he stated the principle plainly: ``What earlier generations held as sacred, remains sacred and great for us too, and it cannot be all of a sudden entirely forbidden or even considered harmful.''\endnote{Benedict XVI, letter to the bishops on the occasion of the publication of \emph{Summorum Pontificum} (7 July 2007), official English text, retrieved and hashed 25 July 2026 (registered artifact; re-fetch byte-identical).} Behind that stands the Council's own limiting norm, which the current \emph{Institutio generalis} itself restates as binding law: ``that innovations not be made in liturgical restoration unless a true and certain benefit of the Church requires it, and with the caution that new forms should in some way grow organically from forms already existing.''\endnote{\emph{Sacrosanctum Concilium} n.~23, restated as a norm binding inculturation at \emph{Institutio Generalis} n.~398, artifact p.~60: ``Norma a Concilio Vaticano II statuta, ut innovationes in instauratione liturgica ne fiant nisi vera et certa utilitas Ecclesiæ id exigat, et adhibita cautela ut novæ formæ ex formis iam exstantibus organice quodammodo crescant \dots{}''} The critical case, at its best, is that this norm is a legal standard with content; that a reform which replaced the offertory prayers wholesale, removed some twenty-four of the twenty-five signs of the cross in the Canon, suppressed Septuagesima, Passiontide and the Pentecost octave, and reduced by an order of magnitude the priest's own prayers cannot easily be described as growth from existing forms; and that the burden of showing ``true and certain benefit'' fell on those who made it and has not obviously been discharged.

The two cases are not symmetrical. The first is an argument about sources, answerable by scholarship; the second is an argument about a standard of prudence, answerable only by judgement---which is why the Church has governed the disagreement rather than adjudicating it.

\subsection{Project synthesis}

\begin{studybox}{Project synthesis: a judgement, its counterargument, and its falsifier}
\textbf{The judgement.} At the loci where the Roman Rite's confession of the Eucharistic sacrifice is legally located---the Canon, the added Eucharistic Prayers, the \emph{Orate, fratres}, the prayer over the offerings, and the instruction that governs them---the postconciliar Missal says what the 1962 Missal says, and in the Roman Canon says it, from \emph{Te igitur} to \emph{Per ipsum}, in the same words. One substantial exception stands inside that sentence and is not softened here: the institution narrative, where three alterations were made---\emph{mysterium fidei} removed from the chalice formula, \emph{Hæc quotiescumque feceritis} replaced by \emph{Hoc facite in meam commemorationem} and moved inside the narrative, \emph{quod pro vobis tradetur} added over the bread---to a form that had stood substantially unchanged in the Latin West for more than a millennium. That is the deepest change in the words of the Roman Canon, and it was made by the promulgating authority, not by drift. It is also a change of a different kind from the rest: what else the reform reduced was not the rite's doctrine but its \emph{redundancy}: the property by which the older book said the same thing many times over, in many registers---repeated prayers, repeated gestures, repeated seasons of anticipation, and a large apparatus of the priest's own speech. The 1962 book named the offering as an immaculate victim at the offertory, again in \emph{Quam oblationem}, again in \emph{Unde et memores}, again in \emph{Supra quæ}, again in \emph{Supplices}, again in \emph{Placeat}; and blessed it with the sign of the cross at each naming. The postconciliar book names it fewer times and blesses it once. On the demonstrable evidence assembled here, the reform's losses are concentrated in that redundancy---in the offertory prayers, the gestural language of the Canon, the priest's private apparatus, the pre-Lenten and Passiontide seasons, and the Pentecost octave---and its gains are concentrated in the word: three readings on Sundays, a three-year cycle, an Old Testament pericope as a normal element, an audible Canon, an obligatory homily, and a set of ministries distributed among distinct books. Both the losses and the gains are real, both were chosen deliberately, and neither is what the standard account of either party says it is.

\smallskip
\textbf{The reasons.}
\begin{itemize}[itemsep=0.15em,topsep=0.25em]
\item \emph{The doctrinal loci were read, not assumed.} \emph{Orate, fratres} with its response and \emph{In spiritu humilitatis} stand unaltered in the newer book (section~4); the Roman Canon stands prayer for prayer from \emph{Te igitur} to \emph{Per ipsum} (section~5); Prayers III and IV carry \emph{Hostia}, \emph{immolatio}, \emph{placari} and \emph{reconciliatio}, which the \emph{Institutio generalis} cites by number as the Missal's own witness to Trent (section~5).
\item \emph{The reductions are countable and are concentrated where redundancy was.} Four offertory prayers dropped against three retained (section~4); twenty-five signs of the cross over the offerings reduced to one (section~5); the celebrant's private prayers reduced to four named moments (\emph{Institutio generalis} n.~33, section~8); Septuagesima, Passiontide and the Pentecost octave suppressed (section~7).
\item \emph{The gains are equally on the page.} Three readings on Sundays and a three-year cycle in place of one, with the Old Testament a normal Sunday element rather than a seasonal one (section~6); an audible Eucharistic Prayer (\emph{Institutio generalis} n.~32, section~8); an obligatory Sunday homily and a universal prayer in the Order itself (section~3); a preface repertory of fifty against fifteen (section~5).
\item \emph{The one substantial textual alteration is named, not absorbed.} The change to the dominical words is stated above and treated at length in section~5, where the 1962 book's own \emph{De defectibus} V contradicts the defence usually offered for it.
\item \emph{Neither book is what its partisans need it to be.} The 1962 volume is itself the product of the 1960 reform and says so on its second leaf (section~3); the current Missal's apparatus cites the 1962 typical edition by name (section~7); and the current law requires vernacular readings in the older rite while directing Latin singing of the Creed in the newer (section~9).
\end{itemize}

\smallskip
\textbf{The strongest counterargument, at full strength.} Redundancy is not decoration; it is how a rite forms the people who use it. A liturgy does not teach by stating a proposition once at the point where a canonist would look for it; it teaches by saturation---by the same truth said in prayer, gesture, silence, posture, season and repetition until it is no longer a proposition but a habit of soul. On that account, a book that says the Mass is a propitiatory sacrifice once in the Canon, once in the \emph{Orate, fratres}, and once in an instruction most laymen will never read has not preserved the doctrine in the sense that matters; it has preserved the doctrine's minimum sufficient statement while dismantling the machinery that transmitted it. It then adds an empirical prediction which is not obviously false: that a generation formed by the reduced book would believe less about the sacrifice than one formed by the fuller book. This study cannot refute that. It can only observe that the argument is about formation rather than text, and that its own judgement is confined to the text.

\textbf{What would overturn the judgement.} Here is the test. Name a proposition of the Council of Trent's session XXII doctrine on the sacrifice of the Mass---that the Mass is a true and proper sacrifice; that it is propitiatory and not merely a sacrifice of praise; that it is offered for the living and the dead, for sins, punishments and satisfactions; that Christ instituted the Apostles priests by the words \emph{Hoc facite in meam commemorationem}; that Masses in honour of the saints are not an imposture---and show that the postconciliar Missal, taken as a whole (the four Eucharistic Prayers, the prayers over the offerings, the prefaces, the Masses for the dead, and the \emph{Institutio generalis}), nowhere expresses it. Any such proposition shown to have no locus in the newer book brings the judgement down, and on the ground that matters most.

The searches this study ran were bounded and returned negative: the propitiatory offering is expressed at Prayer III (\emph{agnoscens Hostiam, cuius voluisti immolatione placari}) and Prayer IV (\emph{sacrificium tibi acceptabile et toti mundo salutare}); the offering for the dead at the Canon's \emph{Memento etiam} and the commemorations in Prayers II--IV; the priesthood conferred at the Supper at \emph{Institutio generalis} nn.~2 and 4, which cite Trent session XXII directly; and the whole is summarised at n.~2, which states that in the new Missal ``the Church's rule of prayer corresponds to her perennial rule of belief,'' that the sacrifice of the cross and its sacramental renewal in the Mass are ``one and the same, save for the manner of offering,'' and that the Mass is therefore ``at once a sacrifice of praise, of thanksgiving, propitiatory and satisfactory.''\endnote{\emph{Institutio Generalis Missalis Romani} (2002), \emph{Proœmium} n.~2, artifact p.~8: ``Ita in novo Missali lex orandi Ecclesiæ respondet perenni legi credendi, qua nempe monemur unum et idem esse, excepta diversa offerendi ratione, crucis sacrificium eiusque in Missa sacramentalem renovationem \dots{} atque proinde Missam simul esse sacrificium laudis, gratiarum actionis, propitiatorium et satisfactorium.''} That search is bounded by the loci examined and correctable; it is a falsifiable claim, not a proof of universal negation.

\smallskip
\textbf{Status.} A source-grounded project judgement by an AI contributor, made from the loci cited and nothing else. It is not magisterial teaching, a canonical opinion or a specialist consensus, and it has had no independent review.
\end{studybox}

\subsection{Three things this study will not say}

It will not say that the two books are the same. A rite whose offertory is four prayers shorter, whose Canon is twenty-four crosses lighter, whose year has three fewer seasonal structures, whose readings triple in volume and whose central prayer is spoken aloud is a different book, and a reader who wants to be told that nothing changed should be told the opposite.

It will not say that the differences are indifferent. Section~4 named a specific and irrecoverable loss at the offertory; section~5 the reduction of a gestural language to a single sign; section~6 the reform's own admission that difficult texts are avoided on Sundays; section~7 three losses in the calendar. Those findings stand whatever anyone concludes from them.

And it will not say which book a reader should prefer, or which a bishop should permit. A comparison of texts cannot reach that question, and the present law reserves the answer to the diocesan bishop and, in the cases the 2023 rescript specifies, to the Apostolic See.

\subsection{A closing observation}

The most durable thing the two books have in common is not a prayer but a phrase. \emph{Ad pristinam sanctorum Patrum normam}: Pius V's revisers claimed it in 1570; the Second Vatican Council used it in 1963; the \emph{Institutio generalis} notes the coincidence and concludes that ``both Roman Missals, although four centuries have intervened, embrace one and the same tradition.'' Whether the second book made good on the claim is what four decades of argument have been about. That both books make it, in the same words, is where any honest comparison has to start.
